\section{A fixed compact support and the full mean reserve}
The initial mean boundary cutoff and the oscillatory packet have
different spatial scales. In
physical coordinates fix two parameters $0<\ell_*\leq1$ and
$0\leq r_*\leq1/4$. Write $c_d$ for the constant in the localized mean
boundary inequality. Let $L_*$ bound all time horizons. The correct
quantity to propagate is
\begin{equation}\label{eq:full-mean-reserve}
 \mathfrak M=\tfrac12 K L_*^2+(B_e+c_d B_c r_*^d)L_*.
\end{equation}
The parameters $K,B_c,B_e$ are nonnegative. Here $K$ bounds the pressure
Hessian above, while $B_c$ and $B_e$ bound
the negative part of the initial symmetric gradient inside and outside
the fixed ball of radius $a_d r_*/\ell_*$, respectively. The boundary
penalty, denoted $\mathcal L$ to distinguish it from $L_*$, is chosen
at least $C_{1,d}B_c$.

\begin{lemma}[Fixed-cutoff propagation]\label{lem:fixed-cutoff}
For a constructed child let $e_p\geq0$ be its upper pressure-Hessian
increment and $e_i\geq0$ its uniform initial-gradient increment. Retain
the same physical $\ell_*,r_*$, replace
\[
 K\mapsto K+e_p,\quad B_e\mapsto B_e+e_i,\quad
 B_c\mapsto B_c+e_i,\quad \mathcal L\mapsto C_{1,d}(B_c+e_i).
\]
All corresponding mean inequalities hold whenever the resulting
$\mathfrak M$ is less than $1/2$, and
\begin{equation}\label{eq:mean-reserve-increment}
 \mathfrak M_{\rm child}-\mathfrak M_{\rm parent}
   =\tfrac12 e_p L_*^2+(1+c_d r_*^d)e_iL_*.
\end{equation}
\end{lemma}
\begin{proof}
The pressure upper bound adds $e_p$. If the difference of the initial
gradients has operator norm at most $e_i$, its quadratic form is at
least $-e_i|z|^2$. The old interior and exterior inequalities therefore
persist on exactly the same regions with the indicated new constants.
The new boundary penalty meets the same coercivity requirement.
Expanding \eqref{eq:full-mean-reserve} proves
\eqref{eq:mean-reserve-increment}. On any shorter horizon the reserve
only decreases. The child's flow, inverse, time derivatives and pressure
are the actual ones already constructed by the packet solve; changing
the parameters for the next mean problem does not replace that child.
\end{proof}

A one-time initialization suffices. If
$b=K L_*^2/2+B_e L_*<1/2$, put
$D=c_dB_c r^dL_*$, $g=1/2-b$, and
$a=\min(1,g/(2(D+1)))$. Replace $\ell,r$ by $a\ell,ar$ once.
The physical splitting radius is unchanged, while
$Da^d\leq Da\leq g/2$, so the full reserve becomes strictly less
than $1/2$. This choice uses only the already constructed initial
stage. Its resulting parameters are frozen for every later stage.

Let $b_*<1/2$ bound that initialized reserve, and use the fixed
dimensional upper bound $\lambda_d=1+c_d4^{-d}$ for $1+c_dr_*^d$.
In the common scale choice, give the primary terms of
$\lambda_d(e_p L_*^2/2+e_iL_*)$ total allowance at most
$(1/2-b_*)/4$, and give their frequency remainders the same total
allowance. The stage estimates split the two errors into these primary
terms and remainders. Summing \eqref{eq:mean-reserve-increment} over
any constructed finite prefix then yields
\[
 \mathfrak M_N\leq\frac{b_*+1/2}{2}<\frac12.
\]
The multiplier $\lambda_d$ depends only on $d$ and is included in the
finite list of construction constants before choosing the starting index
and seed value. The coefficient estimates in
Section~\ref{sec:parameter-envelope} verify the required bounds at the
prescribed frequencies.

\begin{lemma}[One compact set for all initial increments]\label{lem:common-support}
With this fixed-cutoff construction every initial increment is supported
in $\overline B(0,\max(1/4,\ell_*^{-1}))$, irrespective of its
frequency, localization length, or finite truncation order.
\end{lemma}
\begin{proof}
At a stage with physical localization length $0<\rho\leq1$, normalized
labels are $y=x/\rho$, and the mean cutoff is
$\chi((\rho\ell_*)y)$. The actual mean boundary condition expresses
its initial velocity as the compactly supported solenoidal boundary
operator with this cutoff, multiplied by its scalar penalty.
Hence the normalized mean initial field vanishes for
$|y|>(\rho\ell_*)^{-1}$. Physical unscaling replaces $y$ by $x/\rho$,
so it vanishes for $|x|>\ell_*^{-1}$. This is true of every actual
mean grade. The oscillatory profiles and their full tensor correctors
are supported in $|y|\leq1/4$, hence after unscaling in
$|x|\leq\rho/4\leq1/4$. The exact residual-removing correction has
zero initial value. Summing the actual finite family, including its
terminal corrector, proves the assertion. Adding the fixed compact
support of the base datum gives one compact set for all finite-stage
initial velocities and their pointwise limit.
\end{proof}
